\section*{Sources and further reading}
\addcontentsline{toc}{section}{Sources and further reading}
\subsection*{Project sources}
\begin{enumerate}[label={[P\arabic*]},leftmargin=2.5em]
\item Komorebi P. \emph{Projector Refinement Physics, Volume I}. The Abelian gauge and ordered non-Abelian holonomy studies supply the finite field equations and matrix identities. These are written mathematical arguments. Exact formal coverage is identified separately.
\item Komorebi P. \emph{Komorebi Physics formal source}. The gauge verification appendix records the transport and refinement declarations, their hypotheses, and the limitations of the retained compiler evidence.
\item Komorebi P. \emph{Projectors, Packets, and Transfer} and \emph{Fourier Packets and the Continuum Dirac Operator}. The packet framework and free-field operator construction in this book. Extending the latter to a background connection requires the domain and evolution comparisons stated here.
\end{enumerate}

\subsection*{Primary context}
\begin{enumerate}[label={[C\arabic*]},leftmargin=2.5em]
\item C. N. Yang and R. L. Mills. ``Conservation of Isotopic Spin and Isotopic Gauge Invariance.'' \emph{Physical Review} \textbf{96} (1954), 191--195. \url{https://doi.org/10.1103/PhysRev.96.191}.
\item Kenneth G. Wilson. ``Confinement of quarks.'' \emph{Physical Review D} \textbf{10} (1974), 2445. \url{https://doi.org/10.1103/PhysRevD.10.2445}.
\item John Kogut and Leonard Susskind. ``Hamiltonian formulation of Wilson's lattice gauge theories.'' \emph{Physical Review D} \textbf{11} (1975), 395. \url{https://doi.org/10.1103/PhysRevD.11.395}.
\item Pablo Arnault and Fabrice Debbasch. ``Quantum walks and discrete gauge theories.'' \emph{Physical Review A} \textbf{93} (2016), 052301. \url{https://doi.org/10.1103/PhysRevA.93.052301}.
\item C. Cedzich, T. Geib, A. H. Werner, and R. F. Werner. ``Quantum walks in external gauge fields.'' \emph{Journal of Mathematical Physics} \textbf{60} (2019). Author preprint arXiv:1808.10850v1, 31 August 2018. \url{https://arxiv.org/abs/1808.10850v1}; \url{https://doi.org/10.1063/1.5054894}.
\end{enumerate}

\subsection*{How these sources are used}
Yang and Mills provide the original continuum setting for local non-Abelian covariance. Wilson and Kogut--Susskind identify the lattice and Hamiltonian settings in which links, closed-loop observables, electric variables, and constraints belong. These references provide context; their quantization or confinement results are not premises of the elementary identities proved here. A gauge-invariant expression for a walk does not by itself supply a Hamiltonian phase space for its gauge field.

Arnault--Debbasch and Cedzich--Geib--Werner--Werner are prior work on gauge-covariant discrete-time walks. In the latter preprint, equations (19)--(20) give endpoint covariance and conjugation of closed transport, while Section IV.6 explains why coupling depends on the ordered decomposition of a walk. These are direct points of comparison with the conventions and factorwise refinement used here. Their plaquette in equation (22) traverses the opposite orientation to the ratio of forward paths used in this volume; in the Abelian case the two phases are inverses. Its global lattice statements should not be transferred to a periodic quotient without accounting for noncontractible loop holonomy.

The source-check record separates inspected full text from bibliographic authentication. Relevant full-text sections of the Yang--Mills article and the Cedzich et al. preprint were inspected. For Wilson, Kogut--Susskind, and Arnault--Debbasch, the accompanying record authenticated the publisher metadata and abstract, but full-text access through the publisher was unavailable. No uninspected result is used to discharge a proof obligation. This is a focused source comparison, not an exhaustive literature review or a priority claim.
